Definiamo i numeri naturali come segue.

\begin{definition}
    \textbf{Notazione dei numeri naturali}

    L'insieme dei numeri naturali è denotato con $\mathbb{N}$.
\end{definition}

\begin{definition}
    \textbf{Numeri naturali}
    \label{def:naturals}
    
    $ \mathbb{N}$ è definito tramite le due seguenti proprietà caratteristiche.
    \begin{list_2}
        \item $1 \in \mathbb{N}$,
        \item $\forall n \in \mathbb{N}$, $ n + 1 \in \mathbb{N}$ (proprietà del successore).
    \end{list_2}
\end{definition}

\begin{observation}
    I numeri naturali non contengono lo zero, ma è possibile estendere
    l'insieme dei numeri naturali per includere lo zero, definendo così l'insieme $\mathbb{N}_0$.
\end{observation}

\begin{definition}
    L'estensione di $\mathbb{N}$ per includere lo zero è
    \[
        \mathbb{N}_0 := \mathbb{N} \cup \{0\}.
    \]
\end{definition}
    

\vspace{0.7cm}


\subsubsection{Principio di induzione}

Da \textit{\hyperref[def:naturals]{def~\ref*{def:naturals}}} segue direttamente il principio di induzione.

Il principio di induzione è una tecnica dimostrativa tipica delle proprietà delle successioni,
che permette di dimostrare la veridicità di ogni proposizione appartenente a una famiglia
di affermazioni, la quale è in relazione biunivoca con i numeri naturali o un loro sottoinsieme.

\vspace{0.5cm}

\textbf{Principio di induzione:} Sia $\left\{P_n\right\}_{n \geq n_0}$ una famiglia di affermazioni con
$n_0, n \in \mathbb{N}$. Per dimostrare la veridictà di ognuna affermazioni è sufficiente dimostrare le
seguenti due proposizioni.

\begin{list_1}
    \item \textbf{Base dell'induzione:} $P_{n_0}$
    \item \textbf{Passo induttivo:} $\forall n \geq n_0$, $P_n \rightarrow P_{n+1}$
\end{list_1}

\begin{example}
    Si dimostri per induzione 
    \[
        \forall n \in \mathbb{N}, \quad \sum_{k=1}^{n} k = \frac{n(n+1)}{2}.
    \]

    \begin{list_3}
        \item \textbf{Base dell'induzione:}
        \[
            \Hp \;\; n_0=1, \qquad \Th \;\; \sum_{k=1}^{n} k = \frac{(n)(n+1)}{2}
        \]
        \vspace{0.2cm}
        \[
            Dim:
        \]
        \[
            \sum_{k=1}^{1} k = 1 = \frac{1(1+1)}{2}. \qed
        \]
        \item \textbf{Passo induttivo:}
            \[
                \Hp \;\; \sum_{k=1}^{n} k = \frac{n(n+1)}{2}, \qquad \Th \;\; \sum_{k=1}^{n+1} k = \frac{(n+1)(n+2)}{2}
            \]
            \[
                Dim:
            \]
            \[
                \sum_{k=1}^{n+1} k =
                \sum_{k=1}^{n} k + (n+1),
            \]
            Per $Hp$ si ha che
            \[
                \sum_{k=1}^{n} k + (n+1) =
                \frac{n(n+1)}{2} + (n+1),
            \]
            dunque
            \[
                \frac{n(n+1) + 2(n+1)}{2} =
                \frac{(n+1)(n+2)}{2},
            \]
            per proprietà transitiva
            \[
                \sum_{k=1}^{n+1} k =
                \frac{(n+1)(n+2)}{2}. \qed
            \]
    \end{list_3}

\end{example}

\vspace{0.7cm}

\subsubsection{Operazioni in $\mathbb{N}$}

All'interno di $\mathbb{N}_0$ si può definire l'operazione di somma "$+$", che gode delle seguenti proprietà.

\begin{property}
    \textbf{Commutativa}
    \[\forall a, b \in \mathbb{N}_0, \, \, a + b = b + a.\]
\end{property}
\begin{property}
    \textbf{Associativa}
    \[\forall a, b, c \in \mathbb{N}_0, \, \,  (a + b) + c = a + (b + c).\]
\end{property}
\begin{property}
    \textbf{Esistenza dell'elemento neutro}
    \[\forall a \in \mathbb{N}_0, \, \,  a + 0 = 0 + a = a.\]
\end{property}
\begin{property}
    \textbf{Chiusura di $\mathbb{N}_0$ rispetto alla somma} 
    \[\forall a, b \in \mathbb{N}_0, \; (a + b) \in \mathbb{N}_0.\]
\end{property}

\begin{observation}
    All'interno di $\mathbb{N}_0$ non può essere definito l'elemento inverso rispetto alla somma
\end{observation}

